\section{机制讨论}

覆盖收缩主要体现在中纬带，而较高纬带即使覆盖率接近饱和，SWE 仍可发生较大变化（图~\ref{fig:synthesis}）。两种指标在不同纬度带提供互补信息。中纬度可通过覆盖比例识别有雪范围的变化，高纬度的积雪变化在覆盖率接近 1 时更多体现于 SWE。全域原序列中二者共同下降，但去趋势后的低相关和斜率差的不确定性限制了对长期脱耦的判断。

暖冬的低覆盖异常与积雪保存条件变化相容。温度可能关联融化条件和降水相态，现有变量无法估计这些过程的独立贡献。SWE 的当季关联在不同窗口间反号，而全域关联对模型设定敏感。区域内格点异质性、雪变量与温度的共同变化及再分析资料处理均可影响聚合系数。因此即使观测到正关联，降雪输入的变化仍无法辨识（图~\ref{fig:interpretation}）。

\fig{fig11_interpretation}{fig:interpretation}{\textbf{全分析域当季统计关联及待检验过程。} 实线连接 1980--2022 年 43 个 DJF 冬季仅控制线性年份趋势的变量，标注为全域残差回归方向。SWE 关联对 2004 年水平项敏感。虚线连接现有数据无法检验的候选过程。}

早晚期配对位移与年度趋势共同支持分析域阈值边界北移，阈值改变时三个窗口的位移排序发生变化。边界是有雪覆盖南缘的统计量，对经度分布和阈值的变化敏感。靠近掩膜边缘的格点失去连续支撑后，部分经度的边界无法识别。把未识别经度作为缺测，并在共同经度重复分析，可降低覆盖样本变化对纬度趋势的影响。

\section{结论}

\begin{enumerate}
\item 1980--2022 年冬季，ERA5 子集固定分析域内积雪覆盖收缩、近地面气温上升。原序列 SWE 及两种储量估计均下降，加入 2004 年同化变更对应的水平项后，SWE 长期下降的证据减弱。
\item 覆盖率下降在 $30$--$50^\circ$N 最强，原序列 SWE 下降在 $50^\circ$N 以北更快。全域去趋势 SCA 与 SWE 的年际相关接近零，标准化长期趋势差仍未得到明确支持。
\item 控制线性年份趋势后，气温与覆盖率的负关联在中纬度最强。SWE 的温度关联随窗口变化。跨冬季多重校正后，较早冬季气温与随后积雪的关联证据不足。
\item SCF $=0.50$ 的二维地理积雪边界在全分析域中位意义上北移。北移方向在阈值与经度筛选敏感性分析中一致，区域排序和年度趋势证据随定义变化。
\end{enumerate}

\section{局限性}

ERA5 为再分析产品，资料同化随时间变化。已有研究报告 2004 年新增卫星积雪同化可改变雪变量序列的时间一致性~\cite{mortimer2020,urraca2023}。本文的单水平项检验只衡量结果对一次指定变更的敏感程度，无法恢复同质序列，也无法区分同年物理变化与产品变更。窗口和纬度带结果、早晚期边界位移均应在该约束下解读。

SWE 元数据缺少明确的空间代表性定义，高值格点也显著影响储量基线。因此储量只给出两种条件估计。固定掩膜排除了数据包标记的固定冰格点，剩余高 SWE 值的来源无法用三变量辨识。代表窗口采用矩形经纬度界限，边界分析限于有效格点与可识别交叉的经度。

月度输入只覆盖 DJF，限制了对季节内过程和季节外滞后的判断。缺少降雪、降水相态、融雪、升华、辐射与环流变量，无法对覆盖变化或 SWE 变化进行过程归因。积雪边界的纬度指标刻画二维地理位置，海拔变化需要独立地形与积雪资料。
